\documentclass[../Probability.tex]{subfiles}
\begin{document}
\section{Bayesian Regression}
\subsection{Recap}
\subsubsection{Bayes' Rule}
Bayes' Rule (assuming $P(B)>0$):
$$P(A \mid B) = \frac{P(B \mid A)\, P(A)}{P(B)}$$
\begin{itemize}
    \item The prior $P(A)$: Belief about $A$ before seeing $B$.
    \item The likelihood $P(B|A)$: If A true, how probable is the observed evidence $B$.
    \item The marginal likelihood $P(B)$: Total probability of observing $B$ across all scenarios. Acts as normalisation constant.
    \item The posterior $P(A|B)$: Updated belief about $A$ after seeing $B$. 
\end{itemize}
The PDF version:
$$f(x \mid y) = \frac{f(y \mid x)\, f(x)}{f(y)}$$
where $f(y) = \int f(y | x)f(x) \mathrm{d}x$ is the normalising constant.\\
The point is that in Bayesian inference and regression, we will use the form
$$\text{posterior} \propto \text{likelihood} \times \text{prior}$$
ignoring the normalisation constant.

\subsubsection{Unnormalised density functions}
A PDF is defined as:
$$f(x) = \frac{g(x)}{Z}, \quad Z = \int g(x)\, dx$$
$Z$ is the normalising constant. If two functions $g_1$ and $g_2$ differ by only a constant, i.e. $g_1 \propto g_2$, then they define the exact same PDF.\\
$Z$ is often intractable to compute, hence we can just write $f \propto g$ and work with the unnormalised form $g$.

\paragraph{Example} The full PDF of a 1D Gaussian $\mathcal{N}(\mu, \sigma^2)$ is
$$f(x) = \frac{1}{\sqrt{2\pi\sigma^2}} \exp\!\left(-\frac{(x-\mu)^2}{2\sigma^2}\right)$$
Expanding the square in the exponent:
$$f(x) \propto \exp\!\left(-\frac{1}{2\sigma^2}x^2 + \frac{\mu}{\sigma^2}x\right)$$
We dropped terms not involving $x$.\\
Hence if we observe this quadratic form in the exponent, we know it's Gaussian.

\subsubsection{Frequentist vs Bayesian}
Frequentist approach treats Parameters like $\mu$ as fixed, unknown constants, wuth no distributions.\\
MLE is frequentist approach. It finds the parameter that makes the observed data most probable.
\paragraph{Example} A factory produces products labeled 100g. The actual weight $X$ follows a Gaussian distribution with known variance 1, but unknown mean $\mu$.
Three products are sampled with weights 99g, 110g, 107g. What is the best estimate of $\mu$?\\
Given variance = 1 and the independent samples, the likelihood is
$$p(x_i \mid \mu) = \frac{1}{\sqrt{2\pi}} \exp\left(-\frac{(x_i - \mu)^2}{2}\right)$$
The joint likelihood of all observations is (by independence):
$$\mathcal{L}(\mu) = \prod_{i=1}^n p(x_i \mid \mu) = \prod_{i=1}^n \frac{1}{\sqrt{2\pi}} \exp\left(-\frac{(x_i - \mu)^2}{2}\right)$$
Taking the log-likelihood:
$$\log \mathcal{L}(\mu) = -\frac{n}{2}\log(2\pi) - \frac{1}{2}\sum_{i=1}^n (x_i - \mu)^2$$
Maximising this is equivalent to minimising $\sum_{i=1}^n (x_i - \mu)^2$:
$$\frac{d}{d\mu} \sum_{i=1}^n (x_i - \mu)^2 = -2\sum_{i=1}^n (x_i - \mu) = 0$$
$$\boxed{\hat{\mu}_n = \frac{1}{n}\sum_{i=1}^n x_i}$$
MLE says that the sample mean is the best estimate for $\mu$.
However, this approach says nothing about our confidence in this estimate. Also, we cannot incorporate any prior knowledge we may have.\\
\\
Bayesian approach treats parameters as random variables with its own distribution. We can encode our prior beliefs about $\mu$ and then update those beliefs.\\
Treating each observation as Gaussian around $\mu$, the likelihood is:
$$X_i \sim N(\mu, \sigma^2), \quad i = 1, 2, \ldots, n$$
Our prior Gaussian belief about $\mu$ before seeing data:
$$\mu \sim p(\mu) = N(\mu_0, \sigma_0^2)$$
Here $\mu_0$ is our prior best guess for $\mu$, and $\sigma_0^2$ is our uncertainty about that guess.\\
Applying Bayes' Rule,
$$p(\mu \mid x_1, \ldots, x_n) \propto \prod_{i=1}^n p(x_i \mid \mu)\, p(\mu)$$
This is the posterior, i.e. the update belief about $\mu$ after $n$ observations.\\
The posterior is also Gaussian. Its parameters are
$$\boxed{\hat{\mu}_n = \frac{\sigma^2 \sigma_0^2}{n\sigma_0^2 + \sigma^2}\left(\frac{\sum_{i=1}^n x_i}{\sigma^2} + \frac{\mu_0}{\sigma_0^2}\right)}$$
$$\boxed{\hat{\sigma}_n^2 = \frac{\sigma^2 \sigma_0^2}{n\sigma_0^2 + \sigma^2}}$$
The posterior mean can be rewritten as
$$\hat{\mu}_n = \frac{n\sigma_0^2}{n\sigma_0^2 + \sigma^2}\bar{x} + \frac{\sigma^2}{n\sigma_0^2 + \sigma^2}\mu_0$$
where $\bar{x} = \frac{1}{n}\Sigma x_i$ is the MLE estimate.
Note that this is just the weighted average of the MLE estimate and the prior mean. The weights depend on the relative uncertainty.\\
If $\sigma_0^2$ is large, there is more weight on $\bar{x}$. If $\sigma^2$ is large, more weight on $\mu_0$. As $n\to\infty$, $\hat{\mu}_n \to \bar{x}$.\\
The posterior variance $\hat{\sigma}^2_n \to 0$ as $n \to \infty$. Compared to MLE, where $\hat{\sigma}^2_n = 0$ always. 

\subsubsection{Regression}
Recall that Least Squares Method (LSM) finds a linear function of $x$ that best predicts $y$, given $n$ data points, each consisting of input $x_i \in \mathbb{R}^d$ and output $y_i \in \mathbb{R}$.\\
Organising the inputs and outputs:
$$\mathbf{y} = (y_1, \ldots, y_n)^\top \in \mathbb{R}^n$$
$$\mathbf{X}^\top = (x_1, \ldots, x_n) \in \mathbb{R}^{d \times n}$$
We want a weight vector $\mathbf{w} \in \mathbb{R}^d$ such that $\mathbf{Xw} \approx \mathbf{y}$.\\
LSM finds $\mathbf{w}$ by minimising the total squared error:
$$\hat{w} = \arg\min_{w \in \mathbb{R}^d} \|\mathbf{y} - \mathbf{X}w\|^2$$
By differentiating and setting to zero, we get
$$\boxed{\hat{w} = \hat{\Sigma}^{-1}\mathbf{X}^\top\mathbf{y}}$$
where $\hat{\Sigma} = \mathbf{X}^T\mathbf{X}$ and assuming it is invertible. This is the normal equation, the closed-form solution to least squares.\\
If we use MLE by supposing the true relationship between $x$ and $y$ is linear and corrupted by Gaussian noise,
$$y = w^{*\top}x + \varepsilon, \quad \varepsilon \sim N(0, \sigma^2 I)$$
where $w^* \in \mathbb{R}^d$ is the true unknown weight vector.\\
The output $y$ is Gaussian:
$$P(y \mid w, x) = N(w^\top x, \sigma^2 I)(y)$$
By MLE we will also get
$$\hat{w}_{\text{MLE}} = \hat{\Sigma}^{-1}\mathbf{X}^\top\mathbf{y}$$
The implication is that LSM has the same limitation as MLE. It just gives a single estimate without the uncertainty about the prediction. It also doesn't incorporate prior beliefs.
This leads into Bayesian Linear Regression which address these gaps.

\subsection{Bayesian Linear Regression}
\subsubsection{1-D case}
Applying Bayesian inference to the regression setting. Start with 1-D inputs, $x_i\in\mathbb{R}$ and $w\in\mathbb{R}$.\\
\textbf{Likelihood}: Each output is Gaussian around $wx_i$:
$$y_i \sim p(y_i \mid x_i, w) = N(wx_i, \sigma^2), \quad i = 1, 2, \ldots, n$$
\textbf{Prior}: Treat $w$ as a random variable with a Gaussian prior:
$$w \sim p(w) = N(0, \sigma_0^2)$$
The prior is centered at 0, encoding a belief that weights are small.\\
\textbf{The posterior}: By Bayes' Rule,
$$p(w \mid \mathbf{y}) \propto p(\mathbf{y} \mid w, \mathbf{X})\, p(w) = \prod_{i=1}^n p(y_i \mid x_i, w)\, p(w)$$
The result is that the posterior is Gaussian with
\begin{align*}
    \hat{\mu}_n &= \frac{1}{\sigma^2}\left(\frac{1}{\sigma_0^2} + \frac{1}{\sigma^2}\mathbf{X}^\top\mathbf{X}\right)^{-1}\mathbf{X}^\top\mathbf{y}\\
    &= \left(\frac{\sigma^2}{\sigma_0^2}I + \mathbf{X}^\top\mathbf{X}\right)^{-1}\mathbf{X}^\top\mathbf{y}
\end{align*}
$$\hat{\sigma}_n^2 = \left(\frac{1}{\sigma_0^2} + \frac{1}{\sigma^2}\mathbf{X}^\top\mathbf{X}\right)^{-1}$$
% Compared with LSM which gives $\hat{w}_{\text{LSM}} = \left(\mathbf{X}^\top\mathbf{X}\right)^{-1}\mathbf{X}^\top\mathbf{y}$, there is an additional $\frac{1}{\sigma_0^2}$ term, providing the regularisation effect of the prior.
% As $\sigma_0^2 \to \infty$, $\hat{\mu} \to \hat{w}_{\text{LSM}}$.
Comparing with ridge regression which solves $\arg \min_w \|\mathbf{y} - \mathbf{X}w\|^2 + \lambda \|w\|^2$, giving:
$$\hat{w}_{\text{ridge}} = \left(\lambda I + \mathbf{X}^\top\mathbf{X}\right)^{-1}\mathbf{X}^\top\mathbf{y}$$
Note that they are identical with $\lambda = \frac{\sigma^2}{\sigma^2_0}$. So the Gaussion prior $\mathcal{N}(0,\sigma_0^2)$ is mathematically equivalent to $L_2$ regularisation.
\begin{itemize}
    \item Large $\sigma^2$ (noisy data): Large $\lambda$, pulls $w$ towards 0.
    \item Small $\sigma_0^2$ (strong prior): large $\lambda$.
    \item Large $n$, $\mathbf{X}^T\mathbf{X} = \sum_{i=1}^n x_i x_i^\top$ dominates, $\lambda I$ negligible, prior fades away.
\end{itemize}
$\hat{\sigma}^2_n$ represents the remaining uncertainty about $w$ after the observations. It shrinks as $n$ increases.\\
Also, as $\sigma_0^2 \to \infty$ (no prior conviction), $\frac{\sigma^2}{\sigma_0^2} \to 0$ and $\hat{\mu}_n \to \hat{w}_{\text{LSM}}$. Hence LSM is just a special case of Bayesian regression with an infintely weak prior.

\subsubsection{Derivation}
\begin{align*}
p(w|\mathbf{y}) &\propto p(w)p(\mathbf{y}|\mathbf{X}, w) \\
&\propto \exp\left(-\frac{1}{2\sigma_0^2}w^2 - \frac{1}{2\sigma^2}\sum_{i=1}^n(y_i - wx_i)^2\right) \\
&\propto \exp\left(-\frac{1}{2}S_n w^2 + \frac{1}{\sigma^2}\sum_{i=1}^n x_i y_i w\right), 
\qquad S_n = \frac{1}{\sigma_0^2} + \frac{\sum_{i=1}^n x_i^2}{\sigma^2} \\
&\propto \exp\left(-\frac{1}{2}S_n\left(w - \frac{S_n^{-1}}{\sigma^2}\sum_{i=1}^n x_i y_i\right)^2\right) 
\qquad \text{(completing the square)} \\
&\propto N(\hat{\mu}_n, \hat{\sigma}_n)
\end{align*}
$S_n$ collects the $w^2$ coefficients from both the prior and likelihood:
$$S_n = \frac{1}{\sigma_0^2} + \frac{\sum_{i=1}^n x_i^2}{\sigma^2} = \frac{1}{\sigma_0^2} + \frac{\mathbf{X}^\top\mathbf{X}}{\sigma^2}$$
Since a Gaussian has the form $\exp (-\frac{1}{2\hat{\sigma}_n^2}(w - \hat{\mu}_n)^2)$, by matching $S_n = \frac{1}{\hat{\sigma}_n^2}$ and reading off $\hat{\mu}_n$,
$$\hat{\mu}_n = \frac{1}{\sigma^2}\left(\frac{1}{\sigma_0^2} + \frac{1}{\sigma^2}\mathbf{X}^\top\mathbf{X}\right)^{-1}\mathbf{X}^\top\mathbf{y} = \left(\frac{\sigma^2}{\sigma_0^2}I + \mathbf{X}^\top\mathbf{X}\right)^{-1}\mathbf{X}^\top\mathbf{y}$$
$$\hat{\sigma}_n^2 = S_n^{-1} = \left(\frac{1}{\sigma_0^2} + \frac{1}{\sigma^2}\mathbf{X}^\top\mathbf{X}\right)^{-1}$$

\subsubsection{$d$-dimension case}
For $d$ dimensions, scalars just become vectors and scalars.\\
The likelihood and prior become:
$$y_i \sim N(w^\top x_i, \sigma^2 I), \quad w \sim N(\mathbf{0}, \sigma_0^2 I)$$
The posterior is a multivariate Gaussian.
$$\hat{\mu}_n = \frac{1}{\sigma^2}\left(\frac{1}{\sigma_0^2}I + \frac{1}{\sigma^2}\mathbf{X}^\top\mathbf{X}\right)^{-1}\mathbf{X}^\top\mathbf{y}$$
The variance $\hat{\sigma}_n^2$ becomes a covariance matrix $\hat{\Sigma}_n$:
$$\hat{\Sigma}_n = \left(\frac{1}{\sigma_0^2}I + \frac{1}{\sigma^2}\mathbf{X}^\top\mathbf{X}\right)^{-1}$$
Again, as $n \to \infty$, $\frac{1}{\sigma^2}\mathbf{X}^T \mathbf{X}$ dominates and $\hat{\Sigma} \to 0$. The uncertainty drops with more data.

\subsection{Predictive Distribution}
So far Bayesian Linear Regression gives the posterior distribution over weights $p(w | \mathbf{y})$.
But we want to make predictions for new inputs.\\
Given a new input $x$, we want to know the distribution of its output $y$.
This is the predictive distribution $p(y | \mathbf{y}, x)$ (probability of $y$ given all the training data).\\
\\
Here $w$ is a distribution, so every plausible value of $w$ gives a different prediction for $y$. We need to average over all of them, weighted by how probable each $w$ is:
$$p(y \mid \mathbf{y}, x) = \int p(y \mid w, x)\, p(w \mid \mathbf{y})\, dw$$
This is also called marginalising out $w$.\\
Since both $p(y | w,x) = \mathcal{N}(w^\top x, \sigma^2)$ and $p(w | \mathbf{y}) = \mathcal{N}(\hat{\mu}_n, \hat{\Sigma}_n)$ are Gaussian, the result is Gaussian and can be computed as:
$$\boxed{p(y \mid \mathbf{y}, x) = N(\hat{\mu}_n^\top x,\ \sigma^2 + x^\top \hat{\Sigma}_n x)}$$
\begin{itemize}
    \item Predictive mean $\hat{\mu}^\top_n x$: This is simply the posterior mean weights applied to the new input - the same as what LSM would give using $\hat{w}^\top x$. The best single prediction is unchanged.
    \item Predictive variance $\sigma^2 + x^\top \hat{\Sigma}_n x$:
    \begin{itemize}
        \item $\sigma^2$ is the irreducible noise. Even if we knew $w$ exactly, there is inherent randomness in the data generating process $y = w^\top x + \epsilon$.
        \item $x^\top \hat{\Sigma}_n x$ is the model uncertainty. This reflects our uncertainty about $w$ itself, projected onto the new input $x$.
        It depends on where $x$ sits relative to the training data. If $x$ is far from any training point, $\hat{\Sigma}_n$ is large in that direction and this term is large.
    \end{itemize}
\end{itemize}

\subsubsection{Basis Functions}
$x$ can be replaced by basis functions $\phi(x)$. First transform the raw input:
$$x \mapsto \phi(x) = (\phi_1(x), \phi_2(x), \ldots, \phi_k(x))^\top$$
Then the model is still linear in $\phi(x)$ but can still capture nonlinear relationships in the original $x$.

\end{document}